\section*{Chapter \ref{ch:g8:how-the-internet-works} --- How the Internet Works}

\begin{solution}{exo:g8:how-the-internet-works:1}
The four parts, then the 32 \omterm{def:g5:bits-and-bytes:bit}{bits}:
\code{10 -> 00001010}, \code{0 -> 00000000}, \code{0 -> 00000000},
\code{255 -> 11111111}, and
\code{4 bytes, 32 bits: 00001010000000000000000011111111}.
\end{solution}

\begin{solution}{exo:g8:how-the-internet-works:2}
In the order of the numbers 0, 1, 2, 3: \code{"Hel"}, \code{"lo wor"},
\code{"ld"}, \code{"!"}, which make \code{"Hello world!"}.
\end{solution}

\begin{solution}{exo:g8:how-the-internet-works:3}
\code{D -> E -> F in 2 steps}: the \omterm{def:g8:dictionaries:dictionary}{dictionary} gives E after D, and F after E.
\end{solution}

\begin{solution}{exo:g8:how-the-internet-works:4}
The number of pieces is known before cutting: it is the length divided by the
size, rounded up, the trick of the buses of \cref{ch:g6:first-python-program}.
The number of a new piece is one more than the number of pieces made so far:

\omcode{code/grade-8/05-how-the-internet-works/exercises/cut10.py}{1}{12}{Numbered
pieces that know how many there are.}\label{lst:g8:how-the-internet-works:sol-cut}

With \code{"Meet me at the lighthouse at six."} (33 \omterm{def:g7:strings:string}{characters}) and 10,
\code{(33 + 9) // 10} is 4: four pieces, \mbox{\code{"Meet me at"}},
\mbox{\code{" the light"}}, \mbox{\code{"house at s"}} and \code{"ix."}, numbered 1 to 4,
each \mbox{\code{"of": 4}}.
\end{solution}

\begin{solution}{exo:g8:how-the-internet-works:5}
\code{sorted} puts texts in the order of a \omterm{def:g8:dictionaries:dictionary}{dictionary}, \omterm{def:g7:strings:string}{character} by
\omterm{def:g7:strings:string}{character}: \code{"10"} and \code{"11"} start with \code{1}, so they come right
after \code{"1"} and before \code{"2"}. The pieces 10 and 11, \code{k} and
\code{l}, land after \code{b}. Keeping the numbers as numbers sorts them by
\omterm{def:g6:variables-and-input:value}{value}:

\omcode{code/grade-8/05-how-the-internet-works/exercises/rebuild_fixed.py}{1}{8}{Numbers
kept as numbers.}\label{lst:g8:how-the-internet-works:sol-rebuild}
\end{solution}

\begin{solution}{exo:g8:how-the-internet-works:6}
Each part is one \omterm{def:g4:memory-and-storage:byte}{byte}, so it cannot be above 255, and 300 is. The fix adds
that check:

\omcode{code/grade-8/05-how-the-internet-works/exercises/valid_fixed.py}{1}{9}{Each
part from 0 to 255.}\label{lst:g8:how-the-internet-works:sol-valid}

Python checks \code{not part.isdigit()} first, so \code{int(part)} is only
reached when the part is made of digits.
\end{solution}

\begin{solution}{exo:g8:how-the-internet-works:7}
Compare the \omterm{def:g7:lists:list}{lists} of the first three parts:

\omcode{code/grade-8/05-how-the-internet-works/exercises/same_network.py}{1}{3}{The
same first three \omterm{def:g4:memory-and-storage:byte}{bytes}.}\label{lst:g8:how-the-internet-works:sol-same}
\end{solution}

\begin{solution}{exo:g8:how-the-internet-works:8}
\omcode{code/grade-8/05-how-the-internet-works/exercises/sending_time.py}{1}{3}{\omterm{def:g4:memory-and-storage:byte}{Bytes}
to \omterm{def:g5:bits-and-bytes:bit}{bits}, divided by the \omterm{def:g8:how-the-internet-works:bit-rate}{bit rate}.}\label{lst:g8:how-the-internet-works:sol-time}

$5{,}000{,}000 \times 8 = 40{,}000{,}000$ \omterm{def:g5:bits-and-bytes:bit}{bits} at 20,000,000 \omterm{def:g5:bits-and-bytes:bit}{bits} per second:
2 seconds.
\end{solution}

\begin{solution}{exo:g8:how-the-internet-works:9}
$2^{32} = 4{,}294{,}967{,}296$, a little more than four billion.
\end{solution}

\begin{solution}{exo:g8:how-the-internet-works:10}
The speed is $300{,}000 / 1.468 \approx 204{,}360$ kilometres per second, so
$6{,}000 / 204{,}360 \approx 0.029$ seconds: about 29 milliseconds.
\end{solution}

\begin{solution}{exo:g8:how-the-internet-works:11}
Each \omterm{def:g8:how-the-internet-works:packet}{packet} carries $1{,}500 - 40 = 1{,}460$ \omterm{def:g4:memory-and-storage:byte}{bytes} of the \omterm{def:g3:files-and-folders:file}{file}.
$1{,}000{,}000 / 1{,}460 \approx 684.9$, so 685 \omterm{def:g8:how-the-internet-works:packet}{packets}, the last one not
full. They carry $685 \times 40 = 27{,}400$ \omterm{def:g4:memory-and-storage:byte}{bytes} of header, so
$1{,}027{,}400$ \omterm{def:g4:memory-and-storage:byte}{bytes} are sent in all.
\end{solution}

\begin{solution}{exo:g8:how-the-internet-works:12}
A \omterm{def:g8:dictionaries:dictionary}{dictionary} keeps one \omterm{def:g6:variables-and-input:value}{value} per key, so a second copy of a piece simply
replaces the first:

\omcode{code/grade-8/05-how-the-internet-works/exercises/reassemble.py}{1}{13}{Rebuilding
with copies and gaps.}\label{lst:g8:how-the-internet-works:sol-reassemble}
\end{solution}

\begin{solution}{pb:g8:how-the-internet-works:1}
\textbf{1.} $1{,}500 - 40 = 1{,}460$ \omterm{def:g4:memory-and-storage:byte}{bytes}.

\textbf{2.} $3{,}000{,}000 / 1{,}460 \approx 2{,}054.8$: 2,054 full \omterm{def:g8:how-the-internet-works:packet}{packets}
carry $2{,}998{,}840$ \omterm{def:g4:memory-and-storage:byte}{bytes}, and a 2,055th carries the last 1,160 \omterm{def:g4:memory-and-storage:byte}{bytes}. It is
not full.

\textbf{3.} $2{,}055 \times 40 = 82{,}200$ \omterm{def:g4:memory-and-storage:byte}{bytes} of header, out of
$3{,}082{,}200$ \omterm{def:g4:memory-and-storage:byte}{bytes} in all: about 2.7\,\%.

\textbf{4.} The number of \omterm{def:g8:how-the-internet-works:packet}{packets}: the whole number of full \omterm{def:g8:how-the-internet-works:packet}{packets}, plus one
if some \omterm{def:g4:memory-and-storage:byte}{bytes} are left over. Without the \code{if}, the end of the photo would
be lost.

\textbf{5.} 0.098 seconds, about a tenth of a second:

\omoutput{code/grade-8/05-how-the-internet-works/outputs/photo_trip.txt}{The trip of
the photo.}\label{out:g8:how-the-internet-works:photo}

\textbf{6.} 2.466 seconds at 10 megabits per second, 0.247 at 100.

\textbf{7.} The headers travel too: every \omterm{def:g4:memory-and-storage:byte}{byte} sent takes time, not only the
\omterm{def:g4:memory-and-storage:byte}{bytes} of the photo.

\textbf{8.} Sending: 2.466 seconds against 0.098, twenty-five times more.

\textbf{9.} From 10 to 100 megabits per second, the whole trip goes from 2.564
to 0.344 seconds, more than seven times faster. From 100 to 1,000, it only
goes from 0.344 to 0.123 seconds: the sending time shrinks, but the 0.098
seconds of propagation stay.

\textbf{10.} No. Even with an endless \omterm{def:g8:how-the-internet-works:bit-rate}{bit rate}, the first \omterm{def:g5:bits-and-bytes:bit}{bit} needs 0.098
seconds to cover 20,000 kilometres of fibre: nothing travels in a fibre
faster than its light.

\textbf{11.} The \omterm{def:g8:how-the-internet-works:router}{routers} only add time, so the answer stays no; the real trip
is a little longer than 0.098 seconds.

\textbf{12.} \textbf{Result.} The photo travels in 2,055 \omterm{def:g8:how-the-internet-works:packet}{packets}, with 82,200
\omterm{def:g4:memory-and-storage:byte}{bytes} of headers. Crossing 20,000 kilometres of fibre takes 0.098 seconds
whatever the connection, against 2.466 seconds of sending at 10 megabits per
second and 0.247 at 100: past a certain speed, only the distance matters.
\end{solution}
